\chapter{Literature Review}\label{chap:literature}

TBC

\section{Simultaneous Localization and Mapping (SLAM)}\label{sec:literature:slam}

TBC

\subsection{Problem introduction}
\subsection{Scan matching}
\subsection{Extended Kalman Filter (EKF) SLAM}
\subsection{Particle filter}
\subsection{Pose graph}

\section{Systematic Literature Review}\label{sec:literature:systematic}
\todo{TBC}

\subsection{Purpose of the study}\label{sec:literature:systematic:purpose}
The main studies reviewing the \gls{slam} literature are presented in
\tabref{tab:literature:systematic:purpose:current-literature}.
In terms of introducing the problem formulation,
the studies focus on explaining different frameworks for performing \gls{slam}.
Durrant-Whyte~and~Bailey~\cite{durrant-whyte-bailey:ram:2006,bailey-durrant-whyte:ram:2006}
discuss Bayesian-based probabilistic formulations,
namely, \gls{ekf} and the Rao-Blackwellized particle filter frameworks,
also categorized as filtering approaches.
In contrast, Grisetti~\etal~\cite{grisetti:itsm:2010} explain in detail
a smoothing formulation, the graph-based \gls{slam},
which estimates the full trajectory of the robot
from a set of sensor measurements.
Thrun~\cite{thrun:inbook:2008} introduces both probabilistic
(\gls{ekf} and particle filter) and
smoothing (graph) formulations.
Focusing on vision sensorization,
Scaramuzza~and~Fraundorfer~\cite{scaramuzza-fraundorfer:ram:2011,fraundorfer-scaramuzza:ram:2012}
present an extensive tutorial on visual odometry for
estimating relative motion from visual data,
where the study discusses camera modeling and calibration,
motion estimation, and feature matching.
Yousif~\etal~\cite{yousif:iis:2015} extend
the previous visual odometry tutorial
to include methodologies for the vision-based \gls{slam} problem.
Although the introductions mentioned here
provide comprehensive explanations of the problem formulation itself, none of
them focus the discussion on the possible long-term challenges of \gls{slam}.

\begin{table}[t!]
  \caption[Existent Literature Reviews, Surveys, and Tutorials on SLAM]{Existent Literature Reviews, Surveys, and Tutorials on \gls{slam}}
  \label{tab:literature:systematic:purpose:current-literature}
  \centering
  \begin{tabular}{c p{0.43\textwidth} p{0.43\textwidth}}

\hline
\textbf{Year} & \textbf{Topic} & \textbf{Reference}\\
\hline
2006 & Probabilistic formulations (\gls{ekf}, particle filter)
     & Durrant-Whyte~and~Bailey~(2006)~\cite{durrant-whyte-bailey:ram:2006},
       Bailey~and~Durrant-Whyte~(2006)~\cite{bailey-durrant-whyte:ram:2006}\\
\hline
2008 & Probabilistic and pose graph formulations
     & Thrun~(2008)~\cite{thrun:inbook:2008}\\
\hline
2010 & Graph \gls{slam}
     & Grisetti~\etal~(2010)~\cite{grisetti:itsm:2010}\\
\hline
2011 & Observability, convergence, consistency (feature-based \gls{slam})
     & Dissanayake~\etal~(2011)~\cite{dissanayake:iciinfs:2011}\\
\hline
2012 & Visual odometry
     & Scaramuzza~and~Fraundorfer~(2011)~\cite{scaramuzza-fraundorfer:ram:2011},
       Fraundorfer~and~Scaramuzza~(2012)~\cite{fraundorfer-scaramuzza:ram:2012}\\
\hline
2014 & Underwater navigation and localization
     & Paull~\etal~(2014)~\cite{paull:joe:2014}\\
\hline
2015 & Visual \gls{slam}
     & Yousif~\etal~(2015)~\cite{yousif:iis:2015}\\
\hline
2016 & Observability, convergence, consistency (feature and graph-based \gls{slam})
     & Huang~and~Dissanayake~(2016)~\cite{huang-dissanayake:ijars:2016}\\
     & Multi-robot \gls{slam}
     & Saeedi~\etal~(2016)~\cite{saeedi:jfr:2016}\\
     & Visual place recognition
     & Lowry~\etal~(2016)~\cite{lowry:tro:2016}\\
     & \gls{slam} literature overview
     & Cadena~\etal~(2016)~\cite{cadena:tro:2016}\\
\hline
2017 & Autonomous vehicles
     & Bresson~\etal~(2017)~\cite{bresson:tiv:2017}\\
2018 &
     & Kuutti~\etal~(2018)~\cite{kuutti:jiot:2018}\\
\hline
2018 & \gls{ai} for long-term autonomy
     & Kunze~\etal~(2018)~\cite{kunze:ral:2018}\\
     & Long-term sensorization
     & Zaffar~\etal~(2018)~\cite{zaffar:ahs:2018}\\
\hline
2020 & Deep learning
     & Fayyad~\etal~(2020)\cite{fayyad:sensors:2020}\\
     & Multi-robot search and rescue
     & Queralta~\etal~(2020)~\cite{queralta:access:2020}\\
\hline
2021 & Self-driving vehicles
     & Badue~\etal~(2021)~\cite{badue:eswa:2021}\\
\hline
2022 & Underground navigation
     & Ebadi~\etal~(2024)~\cite{ebadi:tro:2024}\\
\hline

  \end{tabular}
\end{table}


Moreover, other studies focus on theoretical aspects
of the \gls{slam} formulation.
Dissanayake~\etal~\cite{dissanayake:iciinfs:2011} discuss the
fundamental properties of \gls{slam}:
observability (whether the problem is solvable),
convergence (whether the state uncertainty converges to a finite value),
and consistency (whether the estimated state is unbiased).
Huang~and~Dissanayake~\cite{huang-dissanayake:ijars:2016}
extend the previous work by providing an in-depth explanation of the
fundamental properties and defining criteria for the
performance evaluation of \gls{slam} algorithms in terms of
consistency, accuracy, and computational efficiency.
In contrast to \cite{dissanayake:iciinfs:2011},
which focuses mainly on filtering-based \gls{slam},
Huang~and~Dissanayake~\cite{huang-dissanayake:ijars:2016} also discuss the
properties for smoothing formulations.
Both studies focus only on theoretical aspects,
without addressing the problem of long-term localization and mapping.

In terms of surveys on \gls{slam} trends,
Cadena~\etal~\cite{cadena:tro:2016} provide a broad overview
of metric and semantic \gls{slam} works.
The work briefly discusses localization and mapping robustness
in terms of loop closure validation and handling dynamic environments,
and \gls{slam} scalability in terms of pose graph sparsification
and parallel and distributed computing.
By contrast, Lowry~\etal~\cite{lowry:tro:2016} focus on topological
\gls{slam}, providing a comprehensive review of visual place recognition.
Although the study discusses the challenges of navigation
across varying conditions, it is limited to vision sensors.
Bresson~\etal~\cite{bresson:tiv:2017} survey trends in single- and
multi-vehicle \gls{slam} and in large-scale experiments for autonomous vehicles.
Although the study compares the methods on accuracy, scalability, availability,
recovery, map updateability, and scene dynamics,
their study only refers to approaches composed of at least odometry and a
mapping module, not discussing localization-only algorithms.
Also, the discussion is more focused on loop closure and relocalization modules
and on leveraging existing data, rather than on methodologies for
addressing the long-term challenges of continuous \gls{slam}.

Similar to \cite{bresson:tiv:2017},
Kuutti~\etal~\cite{kuutti:jiot:2018} and Badue~\etal~\cite{badue:eswa:2021}
analyze trends for self-driving vehicles.
While Kuutti~\etal~\cite{kuutti:jiot:2018} focus on sensorization and
cooperative localization between vehicles,
Badue~\etal~\cite{badue:eswa:2021} concentrate on the architecture of
autonomous driving systems. However, neither of those studies discuss the
challenges of achieving long-term \gls{slam}.
Saeedi~\etal~\cite{saeedi:jfr:2016} present a review of multi-robot \gls{slam}.
Even though the authors identify large-scale, dynamic environments,
multi-session operations, and agent scalability as challenges for
multi-robot \gls{slam}, their study does not provide an overview of
existing methodologies to address these problems.
Paull~\etal~\cite{paull:joe:2014} and Ebadi~\etal~\cite{ebadi:tro:2024}
also present overviews of the \gls{slam} literature,
but are more specific in their domains:
the former focuses on autonomous underwater navigation,
and the latter on \gls{slam} in extreme underground environments.
However, those two studies do not discuss long-term \gls{slam}.

The surveys by Kunze~\etal~\cite{kunze:ral:2018} and
Zaffar~\etal~\cite{zaffar:ahs:2018} focus on
some challenges of long-term autonomy.
Kunze~\etal~\cite{kunze:ral:2018} provides a brief overview of how \gls{ai}
can enable the long-term operation of autonomous systems.
Although the survey discusses challenges such as environments with
varying appearance and the dynamics of moving elements,
the discussion is limited to \gls{ai}-related work.
The work only briefly analyzes localization and mapping tasks,
while also focusing on reasoning, human-robot interaction, and planning.
As for Zaffar~\etal~\cite{zaffar:ahs:2018},
the study focuses solely on reviewing, discussing, and comparing sensors
with respect to sensor lifetime, field operability, ease of replacement,
and suitability for various environments.

Nevertheless, none of the studies presented in
\tabref{tab:literature:systematic:purpose:current-literature} describe
their methodology for selecting the works for discussion.
This limitation prevents other researchers from repeating the review process,
\eg, to update existing studies and maintain an up-to-date understanding
of the current state of the art in \gls{slam}.

As a result, \textit{A systematic literature review on long-term localization
and mapping for mobile robots}~\cite{sousa:jfr:2023} was made in the scope of
this PhD~thesis. The review follows the \gls{prisma}~\cite{page:bmj:2021}
guidelines, defining a systematic methodology to ensure the repeatability of the
selection and data extraction processes while allowing future updates
over the discussion and the review's methodology.
Furthermore, our systematic review does not focus on any specific
strategy or publication time interval.
These considerations improve the study's coverage of
the long-term localization and mapping topic,
compared to the existent \gls{slam} surveys.

In summary, the following questions are the focus of the systematic review on
long-term localization and mapping literature~\cite{sousa:jfr:2023}:
\begin{itemize}
\item main challenges inherent to lifelong \gls{slam};
\item main strategies for accomplishing long-term operations with mobile robots;
\item public datasets commonly used for evaluating long-term localization and
mapping algorithms;
\item how the researchers evaluate the performance of autonomous systems in
long-term operations.
\end{itemize}

\subsection{Review methodology}\label{sec:literature:systematic:method}

PRISMA context:
1. definition of eligibility criteria;
2. conceptualization of the search strategy;
3. selection process:
(a) identification;
(b) screening;
(c) quality assessment.
4. data extraction;
5. synthesis and analysis of the included works.

Eligibility criteria: E1 -- E9 (\tabref{tab:literature:systematic:method:exclusion-criteria})

Search strategy:
- databases
- date and search query
- publication date

Selection process:
- identification
- screening
- quality assessment + graph (\tabref{tab:literature:systematic:method:quality-assessment})

Data extraction process.

Brief analysis number of included records in the review per year!

\begin{figure}[t]
\centering
\includegraphics[height=10cm]{literature/selection.pdf}
\caption[\glsentryfull{prisma} flow diagram for the selection process]{\glslink{prisma}{\glsentryfull{prisma}} flow diagram for the selection process}
\label{sec:literature:systematic:method:prisma-flow}
\end{figure}

\begin{table}[t]
  \caption[Exclusion (E) criteria for the selection process of the systematic literature review]{Exclusion (E) criteria for the selection process of the systematic literature review~\cite{sousa:jfr:2023}.}
  \label{tab:literature:systematic:method:exclusion-criteria}
  \centering
  \begin{tabular}{c m{0.2\textwidth} m{0.68\textwidth}}

\hline
\textbf{E\#} & \textbf{Criteria} & \textbf{Statement}\\
\hline
E1 &
Index &
Papers not indexed in a scientific publication venue\\
\hline
E2 &
Language &
Full-text of the papers not published in English\\
\hline
E3 &
Subject Area &
Papers not classified in the databases as Computer Science, Engineering, Mathematics, or Multidisciplinary\\
\hline
E4 &
Short Papers &
Papers classified as short papers accordingly to the publication venue\\
\hline
E5 &
Gray, Secondary, and Tertiary Literature &
Books, preprints, reports, reviews, thesis, ...\\
\hline
E6 &
Availability &
Full-text of the papers not available in digital libraries\\
\hline
E7 &
Dataset &
Papers that focus only on data collection\\
\hline
E8 &
Coverage &
Papers using only odometry for localization\\
\hline
E9 &
Scope &
Papers that focus on different and not related subjects\\
\hline

  \end{tabular}
\end{table}


\begin{table}[t]
  \caption[Quality Evaluation (QE) criteria and score range for the selection process of the systematic review]{Quality Evaluation (QE) criteria and score range for the selection process of the systematic review~\cite{sousa:jfr:2023}.}
  \label{tab:literature:systematic:method:quality-assessment}
  \centering
  \begin{tabular}{c m{0.74\textwidth} c}

\hline
\textbf{QE\#} & \textbf{Criteria} & \textbf{Score}\\
\hline
QE1 &
Does the paper have an updated state of the art on long-term localization and mapping? &
\{0.0, 0.5, 1.0\}\\
\hline
QE2 &
Is the methodology appropriate and detailed? &
\{0.0, 0.5, 1.0\}\\
\hline
QE3 &
Does the methodology consider both localization and mapping problems? &
\{0.0, 1.0, 2.0\}\\
\hline
QE4 &
Is the hardware and/or software used in the experiments detailed? &
\{0.0, 0.5, 1.0\}\\
\hline
QE5 &
Does the paper presents any kind of long-term experimental results? &
\{0.0, 2.0\}\\
\hline
QE6 &
Does the paper presents comparative results with other methods and/or ground-truth data? &
\{0.0, 1.0\}\\
\hline
QE7 &
Does the work's implementation and/or the data used in the experiments are publicly available? &
\{0.0, 1.0\}\\
\hline
QE8 &
Is the discussion of the results and conclusions appropriate and detailed? &
\{0.0, 0.5, 1.0\}\\
\hline

  \end{tabular}
\end{table}


\subsection{Discussion overview}\label{sec:literature:systematic:discussion}
Refer here that dynamics parts will be discussed later.
Present brief summary on the following topics, based on the categorization made
within the data extraction framework...

\todo{probably, remove the following subsubsections, would make it too long...}

\subsubsection{Appearance variance}
\subsubsection{Map sparsification}
\subsubsection{Multi-session}
\subsubsection{Computational}
\subsubsection{Long-term experimental data}
\subsubsection{Evaluation metrics}

\section{Long-Term SLAM in Dynamic Environments}\label{sec:literature:dynamic}

TBC

\subsection{Map representation}
\subsection{Map matching}
\subsection{Dynamic objects detection}
\subsection{Prediction modeling}

\section{Summary and Conclusions}\label{sec:literature:conclusions}

TBC
